
Five experiments (H-E1, H-M1, H-M2, H-M3, H-M4) address RQ1 through RQ5:

\begin{table}[htbp]
\centering
\resizebox{\columnwidth}{!}{%
\begin{tabular}{llll}
\toprule
Experiment & Research Question & Gate Type & Status \\
\midrule
H-E1 & Existence of benchmark signature & MUST\_WORK & PASS \\
H-M1 & N-gram overlap of removed documents & MUST\_WORK & PASS (dry-run PoC; full experiment pending) \\
H-M2 & Min-k\% memorization differential & SHOULD\_WORK & FAIL (direction reversed at 1B) \\
H-M3 & Contamination-accuracy correlation & SHOULD\_WORK & PASS \\
H-M4 & Token-count vs step-matching robustness & SHOULD\_WORK & PASS \\
\bottomrule
\end{tabular}
}
\end{table}

\subsubsection{Checkpoint Configuration}

\begin{table*}[htbp]
\centering
\resizebox{\dimexpr 2\columnwidth+\columnsep\relax}{!}{%
\begin{tabular}{lllll}
\toprule
Model Size & Pile Step (Primary) & dedup-Pile Step & Token Volume & Token Mismatch \\
\midrule
160M & 99,000 & 143,000 & ~207B & <0.30\% \\
410M & 99,000 & 143,000 & ~207B & <0.30\% \\
1B & 99,000 & 143,000 & ~207B & <0.30\% \\
6.9B & 99,000 & 143,000 & ~207B & <0.30\% \\
\bottomrule
\end{tabular}
}
\end{table*}

H-M2 uses Pile step 98,000 for the 1B PoC (approximately 205.7 billion tokens). H-M4's analytical simulation uses Pile step 128,000 (approximately 218 billion tokens) to represent the step-matched condition.

